% ============================== Part D ==============================
\problem{D1}{Derivatives become multiplications}

\begin{given}
With the Fourier transform $\hat f(\vk) = \int f(\bm\theta)\,e^{i\vk\cdot\bm\theta}\,d^2\theta$ (the sign convention of
B\&S01, Sec.~5.2.1, Eq.~(5.1)), a derivative $\partial/\partial\theta_j$ of $f$ becomes a factor $\mp i k_j$ on
$\hat f$. The sign depends on the convention and will cancel in everything below, because every quantity uses
\emph{two} derivatives. The potential obeys the Poisson equation $\nabla^2\psi = 2\kappa$.
\source{\cBSpoisson.}
\end{given}

\need{C1 ($\kappa$ and $\gamma$ as second derivatives of $\psi$).}

\begin{enumerate}
  \item Show that $\widehat{\partial_i\partial_j f} = -k_ik_j\hat f$ for either sign convention.
  \item Show that $\hat\psi = -2\hat\kappa/k^2$ for $\vk \neq 0$.
  \item What happens at $\vk = 0$? Relate this to C3.
\end{enumerate}

\solution{
(a) Each derivative gives $\mp ik$; two give $(\mp i)^2k_ik_j = -k_ik_j$.
(b) $\nabla^2 \to -(k_1^2 + k_2^2) = -k^2$, so $-k^2\hat\psi = 2\hat\kappa$.
(c) The equation says nothing about $\hat\kappa(0)$, the mean $\kappa$ over the field: shear never constrains it. It is
the mass-sheet degeneracy of C3 in Fourier form.
}

% --------------------------------------------------------------------
\problem{D2}{The Kaiser--Squires kernel}

\begin{given}
The shear is a convolution of $\kappa$ with a kernel. In Fourier space, one multiplication per mode links them, and
the inversion recovers $\hat\kappa$ from $\hat\gamma$ by the complex-conjugate kernel. The papers write the kernel as
$\hat D = \pi(k_1^2 - k_2^2 + 2ik_1k_2)/k^2$ with a factor $\pi^{-1}$ in front of the product; below we absorb the
$\pi$.
\source{\cBSks; \cJeffreyKS.}
\end{given}

\need{C1, D1.}

\begin{enumerate}
  \item Using C1 and D1, show that $\hat\gamma = D\,\hat\kappa$ with $\hat\gamma = \hat\gamma_1 + i\hat\gamma_2$ and
    $D(\vk) = (k_1^2 - k_2^2 + 2ik_1k_2)/k^2$.
  \item Show that $|D| = 1$. Write $D = e^{2i\beta}$, with $\beta$ the direction of $\vk$.
  \item Show that $\hat\kappa = D^*\hat\gamma$ for $\vk \neq 0$. Why is this inversion exact for noise-free data on an
    infinite field, yet noisy in practice?
\end{enumerate}

\solution{
(a) $\hat\gamma_1 = \tfrac12(\widehat{\psi_{11}} - \widehat{\psi_{22}}) = \tfrac12(-k_1^2 + k_2^2)\hat\psi
= \frac{k_1^2 - k_2^2}{k^2}\hat\kappa$ and $\hat\gamma_2 = \widehat{\psi_{12}} = -k_1k_2\hat\psi = \frac{2k_1k_2}{k^2}\hat\kappa$.
(b) $|D|^2 = [(k_1^2 - k_2^2)^2 + 4k_1^2k_2^2]/k^4 = (k_1^2 + k_2^2)^2/k^4 = 1$. With $\vk = k(\cos\beta, \sin\beta)$,
$D = \cos2\beta + i\sin2\beta = e^{2i\beta}$.
(c) $D^*D = |D|^2 = 1$. Noise $\hat n$ in the shear map passes into $\hat\kappa$ undamped ($|D^*| = 1$) at every
scale, including small scales where the signal is tiny. KS therefore needs smoothing afterwards (D4).
}

% --------------------------------------------------------------------
\problem{D3}{Finite fields, padding and the mean mass}

\begin{given}
The inversion holds for $\vk \neq 0$; the real-space form carries an undetermined additive constant $\kappa_0$, because
a uniform surface mass density produces no shear.
\source{\cBSks; \cBSmsd.}
On a grid, the discrete Fourier transform treats the field as periodic.
\end{given}

\need{D1, D2.}

\begin{enumerate}
  \item Why does a periodic transform make the left edge of a shear map act on the right edge? What does padding the
    map with zeros to twice its size do?
  \item Our KS map sets $\hat\kappa(0) = 0$. What does that assume, and can the data test it?
  \item Our likelihood $\kappa$ map pads with zeros too, in the forward direction ($\kappa \to \gamma$). What
    assumption does that make about $\kappa$ outside the field?
\end{enumerate}

\solution{
(a) In a periodic box, mass near one edge is also ``next to'' the opposite edge, and the non-local kernel (shear
falls only as $1/r^2$) couples them. Padding with zeros places the copies far away and removes most of that
wrap-around.
(b) That the mean $\kappa$ in the field is zero. The data cannot test it (D1(c)), so the map's zero point is a
convention.
(c) That $\kappa = 0$ outside the field. Mass outside still shears galaxies inside, so this is an approximation that
matters most near the edge.
}

% --------------------------------------------------------------------
\problem{D4}{\py\ What binning does to a tangential pattern}

\begin{given}
A lens at the origin produces purely tangential shear, $g = -g_t(r)\,e^{2i\phi}$ (B5(b)). KS needs the shear on a
regular grid, so it averages the galaxies in each cell. Smoothing by a Gaussian of width $s$ multiplies each Fourier
mode by $e^{-k^2s^2/2}$ (the convolution theorem). \emph{Derived here}; the cell-averaging argument is ours. Our cells
are 0.98\,Mpc squares and our subhalo sits almost at a cell centre.
\end{given}

\need{B5(b), D2.}

\begin{enumerate}
  \item Show that the average of $g_1$ and $g_2$ over a square centred on the lens is zero for any $g_t(r)$.
    (Use the symmetry of $\cos2\phi$ and $\sin2\phi$ under $\phi \to \phi + \pi/2$ and $\phi \to -\phi$.)
  \item So what does the KS cell containing the subhalo report about the subhalo? Where does KS's information about
    the subhalo come from?
  \item \py\ Compute the cell-averaged shear of an SIS-like subhalo ($g_t = A/r$) for cells whose centre is offset from
    the subhalo by 0 to 1 cell, and compare with the subhalo's shear at the cell centre.
  \item Gaussian smoothing with $s = 0.7$ cell multiplies the mode at the grid's Nyquist scale ($k = \pi/\Delta$, cell
    size $\Delta$) by what factor? What does that mean for the map's resolution?
\end{enumerate}

\solution{
(a) The square is symmetric under a $90^\circ$ rotation, which sends $\cos2\phi \to -\cos2\phi$ and
$\sin2\phi \to -\sin2\phi$ while $r$ is unchanged, so each average equals minus itself.
(b) Nothing (to first order): the cell centred on the subhalo averages its signal away. KS learns about the subhalo
only from the neighbouring cells, where the pattern no longer cancels. The likelihood map, which compares each galaxy
with the shear at its own position, keeps that signal.
(c) See the notebook: the cell average is zero at offset 0 and grows to about the centre value once the cell is
about one cell width away. The nearest cells are where the cancellation hurts.
(d) $e^{-(\pi\cdot0.7)^2/2} = e^{-2.42} = 0.089$. Structure at the scale of one or two cells is suppressed about
tenfold, so KS resolves only features several cells (several Mpc) across.
}
